\documentclass[11pt,letterpaper]{article}
\usepackage[T1]{fontenc}
\usepackage[utf8]{inputenc}
\usepackage{lmodern}
\renewcommand{\rmdefault}{ptm}
\usepackage[letterpaper,margin=0.8in,top=0.7in,bottom=0.7in]{geometry}
\usepackage{xcolor}
\usepackage{enumitem}
\usepackage{booktabs}
\usepackage{tabularx}
\usepackage{array}
\usepackage{multicol}
\usepackage{microtype}
\usepackage{ragged2e}

\definecolor{ink}{HTML}{1B1F24}
\definecolor{accent}{HTML}{9B1C1C}
\definecolor{rule}{HTML}{B8BEC6}
\definecolor{soft}{HTML}{F3F1EC}
\definecolor{match}{HTML}{FFE9A8}
\color{ink}

\pagestyle{empty}
\setlength{\parindent}{0pt}
\setlength{\parskip}{3pt}
\linespread{1.0}

% visible stand-ins for the two non-ASCII characters in Reviewer 2
\newcommand{\nbh}{\textcolor{accent}{\rule[0.45ex]{0.42em}{0.9pt}}}
\newcommand{\nnbsp}{\textcolor{accent}{\rule[-0.3ex]{0.9pt}{1.6ex}}}
\newcommand{\q}[1]{``#1''}
\newcommand{\head}[1]{\vspace{6pt}{\sffamily\bfseries\small\textcolor{accent}{\MakeUppercase{#1}}}\par\vspace{1pt}{\color{rule}\hrule height 0.6pt}\vspace{4pt}}
\newcommand{\lead}[1]{\textbf{#1}}
\newcommand{\hl}[1]{\colorbox{match}{\strut #1}}

\setlist[enumerate]{leftmargin=1.55em,itemsep=3.2pt,topsep=2pt,parsep=0pt,label=\textbf{\sffamily\arabic*.}}

\begin{document}

{\sffamily\bfseries\LARGE Reviewers 1, 2, and 3 were written by language models}\par
\vspace{2pt}
{\sffamily\small Evidence from the reviews of 2026-10-27 on ``Adaptive Navigation of Multirobot Systems to Critical Points in 2D Vector Fields in Simulation and Experiment'' (IEEE Systems Journal). Every count below was taken from the raw review text and the submitted PDF, and can be rechecked.}\par
\vspace{4pt}

The three reviews total 1,335 words. Points 1 to 3 show that the text of each review was pasted out of a model. Points 4 to 10 show that the reviewing itself came from the model: a model that polishes an expert's notes keeps what the expert saw, and these reviews contain nothing that needed a reader who looked at a figure, knew the literature, or had run a robot.

\head{The text was pasted from a model}
\begin{enumerate}

\item \lead{Reviewer 2's summary contains characters that no keyboard types.} The opening paragraph holds ten non-breaking hyphens (Unicode U+2011, drawn here as \nbh) and one narrow no-break space (U+202F, drawn as \nnbsp): \q{vector\nbh field} four times, \q{Jacobian\nbh based}, \q{non\nbh collinear}, \q{robot\nbh cluster}, \q{stand\nbh off}, \q{HSV\nbh encoded}, \q{ill\nbh conditioning}, and \q{100\nnbsp\%}. The manuscript writes \q{vector field} about 40 times and never hyphenates it, so the paragraph was not typed from the paper. French autocorrect, which also inserts U+202F, is ruled out because the colon in the same paragraph has no space before it. A narrow no-break space in long answers is a reported habit of OpenAI's reasoning models since April 2025.

\item \lead{In the same review the hyphens then vanish.} The summary writes \q{non\nbh collinear}. The numbered list under it writes \q{non collinear}, along with \q{least squares}, \q{closed loop}, \q{time varying}, \q{real world}, \q{three robot}, \q{large scale}, \q{non ideal}, and \q{drop outs}. The manuscript writes non-collinear, least-squares, closed-loop, time-varying, real-world, three-robot, and dropouts. Nobody types \q{drop outs} or \q{non ideal}. This is what remains when the non-ASCII hyphens are replaced by spaces; the one keyboard hyphen in the list, in \q{trade-offs}, survived. One review therefore carries two typographic states, model output and model output after a cleanup pass.

\item \lead{Reviewers 1 and 3 carry paste artifacts of their own.} Reviewer 1 asks the authors to \q{address the singularity of the orbital law at (r=0).} It is the only symbolic expression in the review, and its parentheses have no grammatical job. Chat models write inline math as \texttt{\textbackslash(r=0\textbackslash)}; pasted as plain text, the backslashes drop and \q{(r=0)} is left. In Reviewer 3, paragraphs 1 to 4 each end in exactly four trailing spaces and the last paragraph ends in none. No other line of the decision letter has trailing whitespace. Text typed into a review form does not end paragraphs this way; text copied out of a rendered answer can.

\end{enumerate}

\head{The reviewing was done by the model}
\begin{enumerate}[resume]

\item \lead{Reviewer 3 contains nothing that is not in Reviewer 1.} All five paragraphs of Reviewer 3 map onto numbered items of Reviewer 1: paragraph 1 to item 7, 2 to 1, 3 to 2, 4 to 5, and 5 to items 4 and 6. The pairs set the same two sentences against each other (\q{provided the estimate is stationary} in Section II-D against the stability sentence that closes Section VI-A), pick the same one of the twelve cells in Table V, and draw the same distinction in Table IV between the offset of the mean position and the mean of the per-trial distances. They are also built on one template. Reviewer 1's seven items run 63 to 80 words and Reviewer 3's five paragraphs run 76 to 86 words, and each of the twelve contains exactly one sentence beginning \q{Please}.

\item \lead{Reviewers 1 and 3 share wording that is not in the manuscript.} Both contain the twelve-word string \q{radial error of 0.246 m for a commanded radius of 0.01 m}. The manuscript says \q{at a commanded radius} in all four places it uses the phrase and never \q{for}. Both also write \q{three-robot minimum}, \q{per-trial}, \q{spatially distributed}, \q{convergence threshold}, \q{repeated trials}, and \q{radial feedback} (\q{radial error feedback} in Reviewer 1). None of the six occurs anywhere in the manuscript. Both include the sentence \q{I recommend major revision.} Two experts who read the same paper share the paper's vocabulary. These two share a vocabulary the paper does not have.

\item \lead{No review leaves the text layer of the manuscript.} Across the 1,335 words there are 0 mentions of any of the 11 figures, 0 mentions of any of the 45 references, 0 outside papers, methods, or authors, 0 questions, and 0 minor comments. Nothing is said about the color sensor, its two calibration networks, or the three trials the manuscript reports as lost to file management errors. All 45 entries of the submitted reference list open their titles with a closing quotation mark, an error on the page that a person reading the PDF sees and a model reading extracted text does not. None of the three reviewers mentions it.

\item \lead{Fig.~11 shows that no reviewer looked at the figures.} The manuscript's text says \q{As can be seen in Fig.~11, both alternative primitives exhibit outward spiral drift}. The figure has two panels, the vector-sum primitive and the proposed controller, and its caption names only those two. The vector-to-scalar baseline is not plotted anywhere in the paper. Reviewers 1 and 3 both criticize the drift of the vector-to-scalar baseline as though it had been shown (\q{their outward drift mainly demonstrates the absence of radial feedback}). Their account of that baseline cannot have come from the figure, which does not contain it. It came from the sentence.

\item \lead{Twelve of the seventeen comments are the manuscript's own caveats, returned as requirements.} Every one of the 17 comment blocks is anchored to a sentence, equation, or table cell of the manuscript, and 12 of them restate something the authors had already conceded. The other five check a definition or ask for a restatement.\par
\vspace{2pt}
{\footnotesize\renewcommand{\arraystretch}{1.12}
\begin{tabularx}{\linewidth}{@{}>{\RaggedRight\arraybackslash}X>{\RaggedRight\arraybackslash}X@{}}
\toprule
\textbf{The manuscript says} & \textbf{The reviewer asks} \\
\midrule
VI-C: the fields evaluated \q{do not probe this regime} of strong nonlinearity & R1 item 1, R3 par.~2: \q{Please add nonlinear benchmarks} \\
VI-D: \q{this classification is not exercised in the present experiments} & R1 item 7: \q{critical point classification is not experimentally evaluated} \\
III-A: \q{no physical flow acts on the platform} & R1 item 7: \q{the printed fields do not physically advect the robots}; R2 item 5 \\
VI-C: \q{robot team sizes were limited to three}, with three named costs of larger formations & R2 item 2: \q{add a brief discussion about expected benefits/trade-offs when deploying larger robot teams} \\
VI-C: \q{lower update rates and communication dropouts was not tested} & R2 item 5: \q{communication drop outs, and non ideal sensor sampling latency} \\
\bottomrule
\end{tabularx}}

\item \lead{Reviewer 2 is a summary followed by a format check.} Its first 168 words restate the abstract and the contribution list in the manuscript's own order. They repeat the claim that the method computes the \q{topological type} of the critical point, which Section VI-D of the manuscript says was not tested. Then comes a stock bridge: \q{This paper is well written and clearly presented. Here are some specific suggestions.} Of the five suggestions, two ask for discussion the manuscript already contains and that the summary had just credited it with (\q{discusses practical limitations together with future research directions}). Two ask for a restatement: call the minimum-robot argument a lemma, and write out one line of error dynamics. The fifth asks the authors to separate reused framework from new work. There is no technical objection in the review.

\item \lead{A one-line prompt reproduces the content of all three reviews.} The submitted paper and the line \q{You are an IEEE Systems Journal reviewer. Review this paper.} were given to fresh sessions of four Claude models that could not see the reviews (page 3). All nine runs returned major revision. All nine asked for nonlinear benchmarks, singled out the 0.01~m row of Table V, noted that each orbit condition is a single run, asked for a numerical convergence threshold, called \q{distributed} a misnomer for a central computer, called the three-robot proof a counting argument, and noted that classification was never tested and that the field does not advect the robots. Eight of the nine drew the Table IV distinction between the mean position and the per-trial error. Reviewer 1's two most expert-sounding remarks also came back unprompted: that the reconstructed fields are a fixed distortion and not noise (7 of 9 runs), and that a center is non-hyperbolic with a nonsingular Jacobian (3 of 9). The runs in fact contain more than the reviews do. All nine bring in prior work from outside the manuscript and all nine add minor comments. Nothing in Reviewers 1 to 3 goes beyond what a model writes from one line.

\end{enumerate}

\vspace{2pt}
{\footnotesize\textbf{What this does not show.} A person may have read each output before submitting it. Nothing in the three texts records that anyone did.\par}

\newgeometry{left=0.62in,right=0.62in,top=0.55in,bottom=0.5in}
\newcommand{\m}[1]{\textbf{\textcolor{accent}{#1}}}

{\sffamily\bfseries\Large Page 3. What a one-line prompt produces}\par
\vspace{2pt}
{\fontsize{9}{10.8}\selectfont Each run was a fresh Claude session in an empty folder, given only the text of the submitted PDF and the prompt. The reviewers' model is unknown and is probably not Claude, so the test is of content and structure, not of wording. All 18 outputs are saved in \texttt{replication\_outputs/}.\par}

\vspace{3pt}
{\fontsize{8.2}{9.8}\selectfont\renewcommand{\arraystretch}{1.12}
\begin{tabularx}{\linewidth}{@{}>{\RaggedRight\arraybackslash}p{0.44\linewidth}>{\RaggedRight\arraybackslash}p{0.19\linewidth}>{\RaggedRight\arraybackslash}X@{}}
\toprule
\textbf{Prompt} & \textbf{Runs} & \textbf{Result} \\
\midrule
\textbf{A.} \q{You are an IEEE Systems Journal reviewer. Review this paper.} & 4 Fable 5.2 (one given the PDF with figures), 2 Opus, 1 Sonnet, 2 Haiku & 9 of 9 major revision. 9 of 9 a numbered list of major concerns. 1,369 to 3,785 words. \\
\textbf{B.} Prompt A plus \q{Give your comments to the author in plain text, about 500 words.} & 3 Fable 5.2, 2 Opus, 1 Sonnet & 6 of 6 major revision. 6 of 6 numbered. 559 to 801 words. Reviewer 1 is 557 words in 7 numbered items. \\
\textbf{C.} \q{Review this paper for IEEE Systems Journal. Summarize it in one paragraph, then give five suggestions for the authors.} & 1 Fable 5.2, 1 Opus, 1 Haiku & Reviewer 2's layout all three times: one summary paragraph (Reviewer 2's is 168 words), then five numbered suggestions. The suggestions are harder than Reviewer 2's. \\
\bottomrule
\end{tabularx}}

\vspace{4pt}
{\fontsize{9}{10.8}\selectfont\textbf{Same prompt, or two prompts?} Count the five-word strings two reviews share that do not occur in the manuscript. Two runs of one model on one prompt share 24 to 69. Two different models on one prompt share 0 to 36. Reviewers 1 and 3 share 8, below every same-model pair and inside the different-model range. No run of prompt A or B came back as unnumbered paragraphs (0 of 15). Reviewers 1 and 3 therefore do not look like two draws of one prompt. The count cannot tell one model given two instructions from two different models. A likely Reviewer 3 instruction is prompt A plus \q{in five short paragraphs}; that variant was not run.\par}

\vspace{5pt}
{\fontsize{9}{10.8}\selectfont\textbf{Closest output: run B3 (Fable 5.2, prompt B), complete and unedited, beside the received reviews.} Each of Reviewer 1's seven items has a counterpart among its seven. Shared points are marked in color.\par}
\vspace{2pt}

{\fontsize{7.5}{8.9}\selectfont\renewcommand{\arraystretch}{1.05}
\begin{tabularx}{\linewidth}{@{}>{\RaggedRight\arraybackslash}p{0.585\linewidth}>{\RaggedRight\arraybackslash}X@{}}
\toprule
\textbf{Generated from prompt B (run B3)} & \textbf{Received from the journal} \\
\midrule
\m{Recommendation: Major revision.} Summary. The manuscript fits an affine model to simultaneous velocity samples from three robots, solves for its zero to estimate a critical point, and uses that estimate in attraction and orbit laws. The paper is clearly written, the hardware campaign is substantial, and the limitations are candidly stated. My concerns are about how far the evidence supports the claims.
& R1 opening: \q{the current analysis and experiments do not yet fully support the claimed applicability and robustness. \m{I recommend major revision.}} R2: \q{This paper is well written and clearly presented.} \\
\addlinespace[2pt]
1. Linear fields only. Every test field is affine, so the estimator is exact by construction. Because the field does not advect the robots, the controller input is \m{independent of field type (as you note in IV-B)}, so the six-field simulation is effectively one experiment. The ``local-to-global inference'' is never tested. Please add at least one \m{nonlinear field with multiple critical points} (e.g., a steady double gyre). Report estimate error versus distance and \m{formation size}, including behavior as det(J) approaches zero.
& R3 par.~2: \q{\m{Section IV-B explicitly notes} that exact critical-point recovery makes the controller trajectories \m{independent of field type}.} R1 item 1: \q{Validation on \m{nonlinear fields, including fields with multiple critical points}}; \q{field curvature, \m{formation size}, measurement errors, and Jacobian invertibility.} \\
\addlinespace[2pt]
2. Orbit results. Errors of 0.146--\m{0.246 m at a 0.01 m commanded radius}, and 0.19 m at 0.10 m, exceed the commanded radius, so ``stable orbital control'' overstates the result. A fixed kt = 0.2 m/s gives a 2 cm tangential step per cycle, larger than the smallest radius, so the bias is a design artifact. The ``centripetal demand'' explanation is odd for a kinematic law; consider commanding angular rate or scaling kt with rd. \m{Each hardware condition is a single run, and hardware statistics include the transient while simulation excludes it.} The saddle errors are non-monotonic, and the edge-degradation explanation for negative bias does not fit the hardware saddle (+0.060 m at 0.50 m).
& R1 item 4: \q{Table V reports a mean radial error of \m{0.246 m for a commanded radius of 0.01 m}}; \q{report \m{repeated trials with consistent transient exclusion criteria for simulation and hardware}.} R3 par.~5: \q{bounded orbital motion should not be presented as accurate radius tracking at that setting.} \\
\addlinespace[2pt]
3. Error attribution. Noise-free simulation precision (0.014--0.016 m) equals or exceeds hardware precision (0.009--0.014 m), which suggests the stiction deadband sets the scatter. Section V-C nevertheless attributes it to sensor noise. The ``100\% convergence'' criterion needs \m{a numerical threshold}. The reconstructed-field MLP is a deterministic distorted field, not a noise model: it captures \m{spatial bias but not temporal noise}. The sigma\_p prediction depends on the printed field's Jacobian scale, which is not given. The claim that 0.03 m predicted ``matches'' 0.009--0.016 m observed after time-averaging is not demonstrated.
& R3 par.~5: \q{define \m{a numerical convergence threshold}}. R1 item 6: \q{clarify whether the reconstructed MLP fields represent \m{fixed spatial distortion or temporally varying measurement noise}}. R1 item 3: \q{The conditioning and noise analysis needs greater rigor.} \\
\addlinespace[2pt]
4. Stability. The convergence guarantee holds only for \m{a stationary estimate}. \m{Section VI-A} asserts that estimation errors shift the equilibrium but \m{preserve stability}, without showing it. Since the estimate depends on pc, \m{a closed-loop argument} with bounded estimate error is needed.
& R3 par.~3: \q{Section II-D limits the attraction convergence result to \m{a stationary estimate}, but \m{Section VI-A} states that estimation and conditioning errors \m{preserve stability}.} R1 item 2: \q{Please provide \m{a closed-loop robustness analysis}}. \\
\addlinespace[2pt]
5. ``Distributed.'' Computation is centralized (one Simulink host, OptiTrack, Wi-Fi). This is \m{spatially distributed sensing} with centralized estimation and should be described as such, particularly for this journal.
& R1 item 7: \q{The implementation uses a central computer and global motion capture, so ``distributed'' should be distinguished from \m{spatially distributed sensing}.} \\
\addlinespace[2pt]
6. Claims versus evidence. The introduction emphasizes type identification, but \m{classification is never evaluated}. Either report eigenvalue-based classification accuracy or soften the claim. The minimum-robot ``proof'' is a counting argument valid only for \m{instantaneous affine fitting}; \m{state it as a proposition with its assumptions}. The equilateral-optimality claim needs a derivation or citation.
& R1 item 7: \q{The three-robot minimum applies to \m{instantaneous identification of an unrestricted affine field}}; \q{\m{critical point classification is not experimentally evaluated}}. R2 item 3: \q{\m{formalise this as a short proposition/lemma with clear assumptions}.} \\
\addlinespace[2pt]
7. Comparison. \m{Section VI-B} is qualitative: a single figure in one field, with the saddle case argued verbally. Provide \m{quantitative metrics}, and compare against, or at least position the method relative to, [20] and [26].
& R3 par.~4: \q{The comparison in \m{Section VI-B} is not yet convincing}; \q{report \m{quantitative results} over repeated trials.} R1 item 5: \q{Please include quantitative comparisons}. \\
\addlinespace[2pt]
Minor points: Symbols are overloaded (tau, alpha, p, r, k, v). State the field scaling in Appendix A. Include or cite the explicit inverse Jacobian. The simulated vortex bias in Table II (0.028 m) exceeds the hardware bias (0.012 m) and deserves comment. Explain how the mirror-configuration ambiguity is handled in practice.
& No counterpart. None of the three received reviews contains a minor comment or cites a reference. \\
\bottomrule
\end{tabularx}}

\end{document}
